% Generated by scripts/assemble.py; corrections belong in source/corrections.json.
\examhead{83}{哈尔滨工业大学 2026 年数学分析试卷}{81}
% p082-b24; source PDF page 82
\par\Needspace{5\baselineskip}
\blockanchor{p082-b24}
\majorquestion{一}判断题(正确给出理由, 错误给出反例)\par
% p082-b25; source PDF page 82
\blockanchor{p082-b25}
\question{1}函数 $f(x)$ 在开区间 $(a, b)$ 与 $(b, c)$ 均一致连续，则 $f(x)$ 在 $(a, b) \cup (b, c)$ 上一致连续.\par
% p083-b00; source PDF page 83
\blockanchor{p083-b00}
\question{2}函数 $f(x)$ 在 $[a, b]$ 上可导, 且在某点 $c \in (a, b)$ 处 $f'(c) > 0$, 则存在邻域 $(c - \delta, c + \delta)$ 使得 $f(x)$ 在该邻域单调递增.\par
% p083-b01; source PDF page 83
\blockanchor{p083-b01}
\question{3}设函数 $F(x)$ 在 $[a, b]$ 上可导且 $F'(x) = f(x), x \in [a, b]$, 则 $\int_{a}^{b} f(x) \mathrm{d}x = F(b) - F(a)$.\par
% p083-b02; source PDF page 83
\blockanchor{p083-b02}
\question{4}设正项级数 $\sum_{n=1}^{\infty} a_n$ 收敛, 则 $\sum_{n=1}^{\infty} \frac{\sqrt{a_n}}{\sqrt{n}}$ 收敛.\par
% p083-b03; source PDF page 83
\blockanchor{p083-b03}
\question{5}设函数 $f(x, y)$ 在 $(0, 0)$ 处的所有方向极限都存在且相等，则 $\lim_{(x, y) \to (0, 0)} f(x, y)$ 存在.\par
% p083-b04; source PDF page 83
\par\Needspace{5\baselineskip}
\blockanchor{p083-b04}
\majorquestion{二}解答如下问题:\par
% p083-b05; source PDF page 83
\par\Needspace{6\baselineskip}
\blockanchor{p083-b05}
(1) 证明: 如果函数 $f(x)$ 满足: 对任意的 $x \in [a, b]$, 存在 $\delta_x > 0$ 与 $M_x > 0$, 使得\par
\nopagebreak[4]
% p083-b06; source PDF page 83
\blockanchor{p083-b06}
\displayformula{| f (t) | <   M _ {x}, \forall t \in (x - \delta_ {x}, x + \delta_ {x}) \cap [ a, b ].}
% p083-b07; source PDF page 83
\blockanchor{p083-b07}
则 $f(x)$ 在 $[a,b]$ 上有界.\par
% p083-b08; source PDF page 83
\blockanchor{p083-b08}
(2) 如果 (1) 中 $[a, b]$ 换成 $(a, b)$, 结论是否成立? 成立给出证明, 不成立给出反例.\par
% p083-b09; source PDF page 83
\blockanchor{p083-b09}
(3) 如果函数 $f(x)$ 满足在闭区间 $[a, b]$ 上任意一点都存在极限, $f(x)$ 在 $[a, b]$ 上是否有界? 肯定给出证明, 否定给出反例.\par
% p083-b10; source PDF page 83
\par\Needspace{5\baselineskip}
\par\Needspace{6\baselineskip}
\blockanchor{p083-b10}
\majorquestion{三}设 $f(x)$ 在 $x = 0$ 的某个邻域上二阶可微, 且\par
\nopagebreak[4]
% p083-b11; source PDF page 83
\blockanchor{p083-b11}
\displayformula{\lim _ {x \rightarrow 0} \left(1 + x + \frac {f (x)}{x}\right) ^ {\frac {1}{x}} = e ^ {3}.}
% p083-b12; source PDF page 83
\blockanchor{p083-b12}
求 $f(0), f'(0), f''(0)$ 和 $\lim_{x \to 0} \left(1 + \frac{f(x)}{x}\right)^{\frac{1}{x}}$.\par
% p083-b13; source PDF page 83
\par\Needspace{5\baselineskip}
\blockanchor{p083-b13}
\majorquestion{四}设 $f(x)$ 在 $[0,1]$ 上连续可微, 设 $\varepsilon_{n} = \int_{0}^{1} f(x) \mathrm{d}x - \frac{1}{n} \sum_{k=1}^{n} f\left(\frac{k}{n}\right)$, 求极限 $\lim_{n \to \infty} n\varepsilon_{n}$.\par
% p083-b14; source PDF page 83
\par\Needspace{5\baselineskip}
\blockanchor{p083-b14}
\majorquestion{五}设函数 $f(x)$ 在 $(- \infty, + \infty)$ 上可微, 对某点 $a \in (-\infty, + \infty)$, 有 $f(a) = 0$, 且 $|f'(x)| \leq |f(x)|$. 证明: $f(x)$ 在 $(- \infty, + \infty)$ 上恒为 0.\par
% p083-b15; source PDF page 83
\par\Needspace{5\baselineskip}
\blockanchor{p083-b15}
\majorquestion{六}讨论级数  $\frac{1}{1^{p}} - \frac{1}{2^{q}} + \frac{1}{3^{p}} - \frac{1}{4^{q}} + \cdots$  的绝对收敛与条件收敛性.\par
% p083-b16; source PDF page 83
\par\Needspace{5\baselineskip}
\blockanchor{p083-b16}
\majorquestion{七}设函数 $\zeta(x) = \sum_{n=1}^{\infty} \frac{1}{n^x}, x \in (1, +\infty)$.\par
% p083-b17; source PDF page 83
\blockanchor{p083-b17}
(1) 证明: $\zeta(x)$ 在 $(1, +\infty)$ 上连续.\par
% p083-b18; source PDF page 83
\blockanchor{p083-b18}
(2) 证明: $\zeta(x)$ 在 $(1, +\infty)$ 上可微.\par
% p083-b19; source PDF page 83
\blockanchor{p083-b19}
(3) $\zeta(x)$ 在 $(1, +\infty)$ 上是否一致连续？说明理由.\par
% p083-b20; source PDF page 83
\par\Needspace{5\baselineskip}
\blockanchor{p083-b20}
\majorquestion{八}求极限 $\lim_{R\to +\infty}\iint \limits_{|x|\leq R,|y|\leq R}(x^2 +y^2)e^{-(x^2 +y^2)}\mathrm{d}x\mathrm{d}y.$\par
% p083-b21; source PDF page 83
\par\Needspace{5\baselineskip}
\blockanchor{p083-b21}
\majorquestion{九}设 $f(x, y) = \begin{cases} \frac{\ln(1 + xy)}{x}, & x \neq 0; \\ y, & x = 0. \end{cases}$ 证明：$f(x, y)$ 在其定义域上是连续的.\par
% p083-b22; source PDF page 83
\par\Needspace{5\baselineskip}
\par\Needspace{6\baselineskip}
\blockanchor{p083-b22}
\majorquestion{十}计算曲面积分\par
\nopagebreak[4]
% p083-b23; source PDF page 83
\blockanchor{p083-b23}
\displayformula{I = \iint_ {S} \frac {x \mathrm{d} y \mathrm{d} z + y \mathrm{d} z \mathrm{d} x + z \mathrm{d} x \mathrm{d} y}{(2024 x ^ {2} + 2025 y ^ {2} + 2026 z ^ {2}) ^ {\frac {3}{2}}}.}
% p083-b24; source PDF page 83
\blockanchor{p083-b24}
其中 $S$ 是球面 $x^{2} + y^{2} + z^{2} = 1$ 的外侧.\par
